\chapter{Inverse Kinematics}
\label{ch:ik}

\section{From desired pose to joint coordinates}
Inverse kinematics reverses FK. Given a desired grasp position $p_d$ and
orientation $R_d$, find a joint vector $\qvec$ for which
\begin{equation}
 {}^BT_g(\qvec)\approx
 \begin{bmatrix}R_d&p_d\\0&1\end{bmatrix}.
\end{equation}
There may be multiple joint configurations with the same tool pose; some
targets have no solution inside the robot's workspace or joint limits, and
near singular configurations the required joint changes can become very
large. IK is consequently a constrained search, not simply inverting a square
matrix.

\section{The method actually implemented}
The runtime calls SciPy's bounded nonlinear least-squares routine
\code{scipy.optimize.least_squares}. For candidate $q$, it computes the
URDF-based pose and forms a six-element residual:
\begin{equation}
 r(\qvec)=
 \begin{bmatrix}
 p_{FK}(\qvec)-p_d\\
 \mathrm{rotvec}(R_d^TR_{FK}(\qvec))
 \end{bmatrix}\in\R^6.
 \label{eq:ik-residual}
\end{equation}
The first three terms are position error in metres. The last three terms are
the rotation vector (axis times angle) of the relative rotation, in radians.
The implementation concatenates these quantities without scaling or a
separate translation/rotation weighting. Consequently the norm threshold is
a numerical code criterion, not a dimensionally homogeneous physical metric.

The solver bounds are read from URDF limits. Before each solve, the initial
guess is clipped to lie $10^{-8}$ inside each bound, avoiding an initial
coordinate exactly on a boundary. The solve uses at most 150 function
evaluations and sets \code{ftol=xtol=gtol=1e-10}. It accepts a result only if
$\|r(q)\|_2<10^{-6}$; otherwise it tries another seed or raises
\code{ValueError}. The implementation does not provide a custom analytic
Jacobian, collision constraints, posture cost, or task-space damping policy.
SciPy's solver handles its internal local model and numerical derivative
behavior according to the installed SciPy defaults.

\section{Initialization and branch continuity}
For a normal Cartesian sample, the seed is the previous sample's returned
joint vector. This encourages a nearby solution branch and avoids restarting
from a home posture at every millisecond-scale waypoint. The code also rejects
a segment if the maximum adjacent joint step exceeds $0.15$ rad. This is a
discontinuity guard, not a proof that the branch is globally optimal.

When \code{restart=True}, the first guess is followed by 24 additional vectors
sampled uniformly within the URDF limits using a deterministic NumPy random
generator seeded to 7. The task uses restarts for its initial above-pickup
approach IK; each later Cartesian sample is warm-started from its predecessor
without random restarts. Multiple restart attempts improve the chance of
finding a feasible branch but do not enumerate every IK solution.

\section{Worked solver example}
During this guide's source-based numerical check, the runtime solver was asked
for a downward-oriented grasp at
\begin{equation}
 p_d=(0.55,0.35,0.40)\si{m},\qquad R_d=\mathrm{diag}(1,-1,-1),
\end{equation}
with zero seed and deterministic restarts. One returned solution was
\begin{equation}
 \qvec\approx(0.58172582,\ 0.44238155,\ 0.10901428,\ 0,\ -1.23742905,\ -0.58172582)^T\si{rad}.
 \label{eq:ik-example}
\end{equation}
Substituting this vector back into runtime FK gave
$p_{FK}=(0.55,0.35,0.40)$ m to displayed precision and orientation error
$1.9\times10^{-16}$ rad. It is one numerical solution for this solver/version,
not a unique closed-form answer and not a Gazebo task result. A different
valid initial seed may select another joint configuration.

\section{Singularities, unreachable poses, and limits}
A singularity occurs when small tool motion cannot be produced by an ordinary
small change in all independent task directions, or when the inverse mapping
becomes ill-conditioned. The far-pickup optimizer identifies the straight
elbow as its exact horizontal boundary. The selected point is moved inward by
\SI{5}{mm} before route planning to leave margin for tracking and continuous
IK. This is a model-specific constrained, downward-tool envelope, not the
maximum radius for every orientation.

Joint bounds are part of the least-squares problem. A target at $(5,5,5)$ m is
explicitly expected to fail and is covered by an offline test. Other
infeasible goals can fail from nonconvergence even when a simple geometric
reach bound does not rule them out. The segment planner then aborts before
physical motion because the task constructs all route segments before sending
the approach trajectory.

\section{What is not implemented}
The model has a six-axis wrist geometry often associated with closed-form
solutions, but the code does not algebraically solve shoulder/elbow/wrist
branches. It does not build or expose an arm Jacobian, use a pseudoinverse
update $q_{k+1}=q_k+J^+e$, optimize a nullspace objective, or perform
collision-aware IK. Such equations are useful theory, but must not be
presented as this project's method. The actual solver is the residual-based,
bounded nonlinear least-squares procedure described above.

\section{Connection to command generation}
IK returns one $q$ for each Cartesian interpolation sample. The trajectory
code places those joint vectors and estimated velocities into a sequence of
\code{JointTrajectoryPoint} messages. The task node sends the sequence through
the \code{FollowJointTrajectory} action and checks both the action result and
measured joint positions when it completes. Thus successful IK alone does not
move the robot; it produces reference states for the ROS 2 controller.
